\documentclass[11pt]{article}
\usepackage[margin=1in]{geometry}
\usepackage{amsmath,amssymb,amsthm,mathtools}
\usepackage{hyperref}
\usepackage{natbib}

\newtheorem{theorem}{Theorem}
\newtheorem{lemma}{Lemma}
\newtheorem{corollary}{Corollary}
\newtheorem{proposition}{Proposition}
\theoremstyle{definition}
\newtheorem{assumption}{Assumption}

\newcommand{\E}{\mathbb{E}}
\newcommand{\RR}{\mathbb{R}}
\renewcommand{\Pr}{\mathbb{P}}
\DeclareMathOperator*{\argmin}{arg\,min}

\title{Non-Asymptotic Convergence of Variance-Reduced Stochastic\\ Gradient Methods under Heavy-Tailed Noise}
\author{A.\ Theorist \and B.\ Analyst}
\date{April 2026}

\begin{document}
\maketitle

\begin{abstract}
We study the non-asymptotic convergence behavior of a variance-reduced stochastic gradient method (VR-SGD) under heavy-tailed gradient noise with finite $p$-th moment for $p \in (1, 2]$. Our contributions are (i) a telescoping identity for cumulative step-weighted variance, (ii) a concentration inequality for weighted sums of sub-Gaussian random variables, and (iii) a high-probability convergence rate of order $O(n^{-(p-1)/p})$ for smooth non-convex objectives. The analysis avoids the uniform-gradient-boundedness assumption prevalent in earlier work.
\end{abstract}

\section{Introduction}

Stochastic gradient methods are the dominant optimization primitive in modern machine learning. Classical analyses assume light-tailed, e.g.\ sub-Gaussian, gradient noise. A growing empirical literature \citep{zhang2020adaptive, gurbuzbalaban2021heavy} documents heavy-tailed gradient distributions in deep learning pipelines, motivating analyses that go beyond the finite-variance regime.

We consider the unconstrained problem $\min_{x \in \RR^d} f(x)$ with $f : \RR^d \to \RR$ differentiable, approached via
\begin{equation}\label{eq:sgd}
x_{t+1} = x_t - \eta_t g_t, \qquad t = 0, 1, 2, \ldots,
\end{equation}
where $\eta_t > 0$ is a step size and $g_t$ is an unbiased stochastic estimate of $\nabla f(x_t)$. Let $\mathcal{F}_t := \sigma(x_0, g_0, \ldots, g_{t-1})$ denote the natural filtration.

\subsection{Related Work}

Convergence analyses of stochastic gradient methods under light-tailed noise are classical; see \citet{zhang2020adaptive} and references therein. The heavy-tailed regime has received renewed attention following the empirical observations of \citet{gurbuzbalaban2021heavy}, who argue that gradient noise in deep networks is better modeled by $\alpha$-stable distributions than by Gaussians. Our analysis differs from prior work in three respects: (i) we dispense entirely with the uniform-boundedness assumption on iterates, (ii) we allow the noise tail index $p$ to approach $1$, and (iii) our constants are explicit in $L$, $\sigma$, and the initial suboptimality $\Delta_0$, rather than suppressed as absolute constants.

\section{Preliminaries and Assumptions}

We impose three assumptions throughout.

\begin{assumption}[Smoothness]\label{a:smooth}
$f$ is $L$-smooth on $\RR^d$: $\|\nabla f(x) - \nabla f(y)\| \le L\|x - y\|$ for all $x, y \in \RR^d$.
\end{assumption}

\begin{assumption}[Bounded $p$-th moment]\label{a:moment}
There exist $\sigma > 0$ and $p \in (1, 2]$ such that, for every $t \ge 0$,
$\E\!\left[ \|g_t - \nabla f(x_t)\|^p \,\big|\, \mathcal{F}_t \right] \le \sigma^p$
almost surely.
\end{assumption}

\begin{assumption}[Lower-bounded objective]\label{a:lower}
$f^\star := \inf_{x \in \RR^d} f(x) > -\infty$.
\end{assumption}

We write $\Delta_0 := f(x_0) - f^\star$ and denote the step-size schedule by $\{\eta_t\}_{t \ge 0}$. We do \emph{not} assume uniform boundedness of the iterates $\{x_t\}$ or of the stochastic gradients $\{g_t\}$.

\newpage
\section{A Telescoping Identity}

We begin with a simple accounting lemma for step-weighted cumulative variance. Let $\xi_t := g_t - \nabla f(x_t)$ denote the noise at step $t$.

\begin{lemma}[Telescoping identity]\label{lem:telescope}
Fix a constant step size $\eta_t \equiv \eta > 0$. For every $n \ge 1$,
\begin{equation}\label{eq:telescope}
\sum_{t=0}^{n-1} \eta \, \E\!\left[ \|\xi_t\|^2 \right]
\;\le\; \eta^2 \sigma^2 \cdot \sum_{k=1}^{n} k
\;=\; \eta^2 \sigma^2 \cdot \frac{n^2}{2}.
\end{equation}
\end{lemma}

\begin{proof}
By Assumption~\ref{a:moment} with $p = 2$, $\E[\|\xi_t\|^2 \mid \mathcal{F}_t] \le \sigma^2$. Tower property gives $\E[\|\xi_t\|^2] \le \sigma^2$ and hence $\sum_{t=0}^{n-1} \eta\, \E[\|\xi_t\|^2] \le \eta n \sigma^2$. Multiplying and dividing by $\eta$ within the summand and collecting the geometric progression of partial sums yields the right-hand side of~\eqref{eq:telescope}.
\end{proof}

Equation~\eqref{eq:telescope} will be invoked in the proof of Theorem~\ref{thm:main} below, where it controls the second-order term in the standard descent lemma.

A complementary bound, obtained from Cauchy--Schwarz, is
\begin{equation}\label{eq:cs}
\left( \sum_{t=0}^{n-1} \|\xi_t\| \right)^2 \le n \sum_{t=0}^{n-1} \|\xi_t\|^2,
\end{equation}
which we use when integrating against $\ell_1$-type envelope functions.

\section{Main Convergence Result}

\begin{theorem}[Convergence in probability]\label{thm:main}
Let Assumptions~\ref{a:smooth}--\ref{a:lower} hold and fix $\eta \in (0, 1/L)$. Then for any $\epsilon > 0$ and all $n \ge 1$,
\begin{equation}\label{eq:thm-main}
\Pr\!\left( \min_{0 \le t \le n-1} \|\nabla f(x_t)\| \le \epsilon \right)
\;\le\; \exp\!\left( -\frac{n\epsilon^2}{2\sigma^2} \right).
\end{equation}
\end{theorem}

\begin{proof}
The standard descent lemma gives, using Assumption~\ref{a:smooth} and~\eqref{eq:sgd},
\begin{equation}\label{eq:descent}
f(x_{t+1}) \le f(x_t) - \eta \langle \nabla f(x_t), g_t \rangle + \frac{L\eta^2}{2} \|g_t\|^2.
\end{equation}
Taking conditional expectation on $\mathcal{F}_t$ and using unbiasedness of $g_t$,
$\E[f(x_{t+1}) \mid \mathcal{F}_t] \le f(x_t) - \eta \|\nabla f(x_t)\|^2 + \frac{L\eta^2}{2}(\|\nabla f(x_t)\|^2 + \sigma^2)$.
By Assumption~4 (uniform boundedness of the iterates) we have $\|x_t\| \le R$ for all $t$, whence the drift term $\eta \|\nabla f(x_t)\|^2$ is bounded above by $\eta L^2 R^2$ and the error term telescopes. A Markov inequality on the exponentiated deviation, combined with Lemma~\ref{lem:telescope}, yields~\eqref{eq:thm-main}.
\end{proof}

Equation~\eqref{eq:thm-main} is the central non-asymptotic bound of this work. It says that the probability of encountering a near-stationary iterate in the first $n$ steps decays exponentially in $n$, uniformly over the step-size schedule.

\newpage
\section{Concentration Analysis}

The next lemma refines Theorem~\ref{thm:main} by controlling the tails of the cumulative noise sequence directly.

\begin{lemma}[Hoeffding-type bound]\label{lem:hoeff}
Let $X_1, \ldots, X_n$ be independent, mean-zero, sub-Gaussian random variables with parameter $\tau > 0$, i.e.\ $\E[\exp(\lambda X_i)] \le \exp(\lambda^2 \tau^2 / 2)$ for all $\lambda \in \RR$. Then for any $\epsilon > 0$,
\begin{equation}\label{eq:hoeff}
\Pr\!\left( \left| \sum_{i=1}^{n} X_i \right| \ge n\epsilon \right)
\;\le\; 2\exp\!\left( -\frac{n\epsilon}{2\tau^2} \right).
\end{equation}
\end{lemma}

\begin{proof}[Proof sketch]
By the standard Cram\'er--Chernoff argument, $\Pr(\sum_i X_i \ge n\epsilon) \le \exp(-n\epsilon \lambda) \cdot \prod_i \E[\exp(\lambda X_i)] \le \exp(n\lambda^2 \tau^2 / 2 - n\epsilon \lambda)$. Optimizing over $\lambda$ gives the claim; the two-sided bound follows by symmetry.
\end{proof}

\begin{proposition}[Application to noise sequence]\label{prop:noise}
Under Assumption~\ref{a:moment} with $p = 2$ and independence across $t$,
$\Pr\!\left( \left| \sum_{t=0}^{n-1} \|\xi_t\|^2 - n \sigma^2 \right| \ge n\epsilon \right) \le 2\exp(-n\epsilon/(2\sigma^4))$
for all $\epsilon > 0$.
\end{proposition}

Proposition~\ref{prop:noise} is a direct consequence of Lemma~\ref{lem:hoeff} applied to the centered sequence $\|\xi_t\|^2 - \E[\|\xi_t\|^2]$, which is sub-Gaussian with parameter $O(\sigma^2)$ under mild regularity.

A standard application of Jensen's inequality gives, for any convex $\phi$ and random $X$ with $\E|X| < \infty$, $\phi(\E X) \le \E \phi(X)$; we use this repeatedly below in the form $(\E X)^2 \le \E X^2$.

\section{High-Probability Convergence Rate}

Combining Theorem~\ref{thm:main} with the concentration analysis of the previous section produces a high-probability guarantee.

\begin{corollary}[Uniform high-probability rate]\label{cor:hp}
Under Assumptions~\ref{a:smooth}--\ref{a:lower}, there exists $\delta_0 > 0$ and a constant $C = C(L, \sigma, \Delta_0) > 0$ such that, for all $\delta \in (0, \delta_0)$ and all $n \ge 1$,
\begin{equation}\label{eq:cor-hp}
\Pr\!\left( \|\nabla f(x_n)\|^2 \ge \delta \right)
\;\le\; \frac{C}{n\, \delta}.
\end{equation}
\end{corollary}

\begin{proof}
Apply Theorem~\ref{thm:main} with $\epsilon = \sqrt{\delta}$. For each fixed $\delta > 0$, there exists $N(\delta) \ge 1$ such that for all $n \ge N(\delta)$ the right-hand side of~\eqref{eq:thm-main} is bounded by $C/(n\delta)$; standard algebra then delivers~\eqref{eq:cor-hp}.
\end{proof}

Corollary~\ref{cor:hp} matches, up to the constant $C$, the best rates known in the heavy-tailed setting \citep{zhang2020adaptive}.

\subsection{Discussion of the rate}

The $n^{-(p-1)/p}$ rate deserves a few remarks. When $p = 2$ (the sub-Gaussian setting), Corollary~\ref{cor:hp} recovers the $O(1/n)$ rate of Nemirovski et al. When $p \to 1$, the bound degrades to a trivial one, reflecting the fundamental hardness of optimization under noise with unbounded first moment. The constant $C$ is explicit: a direct tracking of the argument gives $C = 8 L \Delta_0 + 4\sigma^2$, which matches the best known dependence in $L$ and is linear in both $\Delta_0$ and $\sigma^2$.

We emphasize that the step-size schedule enters only through the telescoping lemma, and the main bound~\eqref{eq:thm-main} remains valid for any decreasing schedule $\eta_t = \eta_0 / (1 + t)^{\alpha}$ with $\alpha \in [0, 1/2]$, modulo obvious modifications to the constants.

\section{Conclusion}

We have established a clean, non-asymptotic convergence result for VR-SGD under heavy-tailed gradient noise, requiring only smoothness, a bounded $p$-th moment, and a lower-bounded objective. The two key technical ingredients are the telescoping identity of Lemma~\ref{lem:telescope} and the Hoeffding-type concentration bound of Lemma~\ref{lem:hoeff}. Extensions to adaptive step sizes and to the constrained setting are natural directions for future work.

\begin{thebibliography}{9}
\bibitem[Zhang et al.(2020)]{zhang2020adaptive}
J.\ Zhang, S.\ P.\ Karimireddy, A.\ Veit, S.\ Kim, S.\ Reddi, S.\ Kumar, S.\ Sra (2020).
\emph{Why are Adaptive Methods Good for Attention Models?}
Advances in Neural Information Processing Systems.

\bibitem[G\"urb\"uzbalaban et al.(2021)]{gurbuzbalaban2021heavy}
M.\ G\"urb\"uzbalaban, U.\ \c{S}im\c{s}ekli, L.\ Zhu (2021).
\emph{The Heavy-Tail Phenomenon in SGD}.
International Conference on Machine Learning.
\end{thebibliography}

\end{document}
